\section{Related work}
\label{sec:related}

\paragraph{Spatio-temporal traffic forecasting}
Graph-based recurrent and convolutional models \citep{li2018dcrnn,yu2018stgcn,wu2019graphwavenet,bai2020agcrn} couple temporal sequence modelling with spatial message passing over a road or learned graph. Attention-based architectures \citep{vaswani2017attention} have since been adapted to spatio-temporal data; ReShift uses STLAformer \citep{stlaformer} as its backbone but does not depend on its internals beyond a latent representation and the forecast. Classical non-parametric forecasting based on nearest historical neighbours \citep{davis1991nonparametric,smith1997traffic} is an early form of retrieval: future traffic is predicted from the futures of similar past states. ReShift revisits this idea inside a deep forecaster and under missing inputs.

\paragraph{Forecasting with missing observations}
Imputation methods reconstruct missing entries using recurrent dynamics \citep{che2018grud,cao2018brits}, graph neural networks \citep{cini2022grin}, low-rank tensor or matrix structure \citep{chen2019bayesian} or conditional diffusion models \citep{tashiro2021csdi}. Their reconstructions can be passed to any forecaster, but reconstruction error propagates and the imputer is not trained for the downstream forecasting loss. Direct approaches make the forecaster aware of the mask, for example through decay mechanisms \citep{che2018grud}, joint local and global modelling \citep{tang2020joint} or masking-aware recurrent architectures for network-wide traffic \citep{cui2020missing}. ReShift is complementary: it starts from such a missing-aware forecaster and estimates what remains of the failure-induced error in output space.

\paragraph{Retrieval-augmented prediction}
Retrieval augmentation conditions a parametric model on examples fetched from a datastore, as in nearest-neighbour language models \citep{khandelwal2020knnlm} and retrieval-augmented generation \citep{lewis2020rag}. For time series, retrieval is attractive because recurrent daily and weekly patterns make the training set an informative memory. Two difficulties are specific to the incomplete setting: queries and memory items are described under different observation conditions, which we address with contrastive key alignment \citep{oord2018cpc} and a momentum memory encoder \citep{he2020moco}; and retrieved futures may be wrong for the query, which we address with per-element reranking against a no-correction option and risk gating.

\paragraph{Counterfactual targets and distillation}
Using a fixed model to produce training targets for another model is the essence of knowledge distillation \citep{hinton2015distilling}. ReShift uses a frozen copy of the pretrained forecaster as a teacher, but distils a quantity that is not a prediction of the data: the teacher's own counterfactual change of output when the missing entries are restored. We are not aware of work that treats this restoration shift as the object of estimation and estimates it jointly from retrieved, replayed and parametric evidence for traffic forecasting \todo{confirm with a systematic literature search}.
